\section{Finite source identities and decomposition}
\label{sec:source-decomposition}

Throughout this section, \(N\) is an even integer with \(N\ge2\).
Further size hypotheses appear at each display or just before it.
Use \(L=\log N\) from Section~\ref{sec:notation}.
We prove the finite identities of this section below.
They count each occurrence with its multiplicity.
The later analytic estimates are separate hypotheses.
Use \(\mathcal A_N,G_b(N),\varepsilon,a,b\) from
\eqref{eq:intro-exponents}--\eqref{eq:intro-count} and
\(\alpha,\beta,\mu_*,\lambda,t_0,u\) from
\eqref{eq:not-sieve-parameters}. Take the roughness cutoff
\(z_s=N^\alpha\) from the Root--A scale definition in
\eqref{eq:ra-parameters}.
Throughout, \(n\) is a positive source integer. When
\(n\in\mathcal A_N\), \(n=N-p\) for a prime \(p\). In a source sum, \(e\)
is the integer being counted. The letters \(p,p_i,a_i,r\) denote
prime variables, \(q\) is a positive modulus, and \(m\) is the factor or
cofactor specified by its sum. Paper-II counts use
\(n=dm\); the marked count of Section~\ref{sec:root-a} uses \(m=(N-e)/p\). These roles
are local to the displayed source.
Here \(\Omega(n)\) counts prime factors with multiplicity.
For squarefree \(d\) with \((d,N)=1\), define
\begin{equation}\label{eq:source-paperii-cofactor}
 S^{\mathrm{II}}(d;y)
 =\#\{n=dm\in\mathcal A_N:P^-(m)\ge y,\ (m,Nd)=1\}.
\qquad(2\mid N,\ N\ge2).
\end{equation}
Set this count to zero for other \(d\).
For admissible \(d\), it is Paper~II's (1.7) count in our notation
\cite[(1.7)]{PaperII}. The zero value for other \(d\) is our convention.
Put \(x_i=\log p_i/\log N\).
The tenth row uses
\begin{equation}\label{eq:source-paperii-clipping}
\begin{aligned}
 C_{10}&=\frac{57}{64}-\frac1{500}-\frac{11}{5}\alpha,\\
 B_{10}(x)&=\begin{cases}
 \beta,&\beta\le x<\mu_*,\\
 \min\{\beta,\max\{\alpha,C_{10}-5x/2\}\},
     &\mu_*\le x<\lambda.
 \end{cases}
\end{aligned}
\qquad(2\mid N,\ N\ge2).
\end{equation}
Table~\ref{tab:source-paperii-rows} lists all thirteen sums with every
divisor and cutoff of \cite[(1.6)--(1.7), row table]{PaperII}.
All \(p_i\) in the table are prime.
Each admissible divisor contributes separately.

\begin{longtable}{@{}c c p{.47\textwidth} p{.19\textwidth}@{}}
\caption{The Paper~II source sums. Each row sums \(S^{\mathrm{II}}(d;y)\).}
\label{tab:source-paperii-rows}\\
\hline
Row & Coefficient & Divisor and domain & Cutoff \(y\)\\
\hline
\endfirsthead
\hline
Row & Coefficient & Divisor and domain & Cutoff \(y\)\\
\hline
\endhead
\(S_1\)&\(3\)&\(d=1\)&\(N^\alpha\)\\
\(S_2\)&\(1\)&\(d=1\)&\(N^\beta\)\\
\(S_3\)&\(-2\)&\(d=p_1,\ \alpha\le x_1<\mu_*\)&\(N^\alpha\)\\
\(S_4\)&\(-1\)&\(d=p_1,\ \mu_*\le x_1<1/4\)&\(N^\alpha\)\\
\(S_5\)&\(-1\)&\(d=p_1,\ 1/4\le x_1<t_0\)&\(N^\alpha\)\\
\(S_6\)&\(-1\)&\(d=p_1,\ t_0\le x_1<1/3\)&\(N^\alpha\)\\
\(S_7\)&\(-1\)&\(d=p_1,\ \mu_*\le x_1<\lambda\)&\(N^\alpha\)\\
\(S_8\)&\(1\)&\(d=p_1p_2,\)\newline
 \(\alpha\le x_2<x_1<\beta\)&\(N^\alpha\)\\
\(S_9\)&\(1\)&\(d=p_1p_2,\)\newline
 \(\alpha\le x_2<\beta\le x_1<\mu_*\)&\(N^\alpha\)\\
\(S_{10}\)&\(1\)&\(d=p_1p_2,\ \mu_*\le x_1<\lambda,\)\newline
 \(\alpha\le x_2<B_{10}(x_1)\)&\(N^\alpha\)\\
\(S_{11}\)&\(-2\)&\(d=p_1p_2,\ \lambda\le x_1<1/3,\)\newline
 \(x_1<x_2<(1-x_1)/2\)&\(p_2\)\\
\(S_{12}\)&\(-1\)&\(d=p_1p_2,\)\newline
 \(\alpha\le x_1<1/3\le x_2<(1-x_1)/2\)&\(p_2\)\\
\(S_{13}\)&\(-1\)&\(d=p_1p_2,\)\newline
 \(\beta\le x_1<\lambda\le x_2<(1-x_1)/2\)&
 \(\sqrt{N/(p_1p_2)}\)\\
\hline
\end{longtable}

Let \(\mathcal D_i\) be the displayed divisor domain and \(y_i(d)\) its cutoff.
Define the row occurrences by
\begin{equation}\label{eq:source-row-occurrences}
 \vartheta_i(n)=\sum_{\substack{d\in\mathcal D_i\\d\mid n\\
                  d\ \mathrm{squarefree},\ (d,N)=1}}
  1_{P^-(n/d)\ge y_i(d)}1_{(n/d,Nd)=1}.
\qquad(2\mid N,\ N\ge2).
\end{equation}
Thus \(S_i=\sum_{d\in\mathcal D_i}S^{\mathrm{II}}(d;y_i(d))\).
Keep every cofactor and endpoint restriction. Define
\begin{equation}\label{eq:source-thirteen}
\begin{aligned}
 S_i&=\sum_{n\in\mathcal A_N}\vartheta_i(n)\qquad(1\le i\le13),\\
 W_{13}(n)&=3\vartheta_1(n)+\vartheta_2(n)-2\vartheta_3(n)
 -\vartheta_4(n)-\vartheta_5(n)-\vartheta_6(n)-\vartheta_7(n)\\
 &{}+\vartheta_8(n)+\vartheta_9(n)+\vartheta_{10}(n)
 -2\vartheta_{11}(n)-\vartheta_{12}(n)-\vartheta_{13}(n),\\
 s_{\mathrm{sw}}(n)&=1_{\Omega(n)\ge3}\max\{W_{13}(n),0\},\\
 S_{\mathrm{sw}}&=\sum_{n\in\mathcal A_N}s_{\mathrm{sw}}(n),\qquad
 w(n)=W_{13}(n)-s_{\mathrm{sw}}(n),\qquad
 \mathcal W=\sum_{n\in\mathcal A_N}w(n).
\end{aligned}
\qquad(2\mid N,\ N\ge2).
\end{equation}
This is the complete switched source of Paper~II, (1.9)
\cite[(1.9)]{PaperII}.
Every divisor prime and cofactor cutoff is at least $N^\alpha$.
For $S_{13}$ the cutoff exceeds $N^{(1-\lambda)/4}>N^\alpha$.
Thus every $\vartheta_i(n)$ vanishes when $P^-(n)<N^\alpha$,
so the extra roughness indicator of Paper~II is redundant.
Use $P^-(1)=\infty$ as in Table~\ref{tab:notation}.
Finite linearity gives
\begin{equation}\label{eq:source-weight-identity}
\begin{aligned}
 \mathcal W={}&3S_1+S_2-2S_3-S_4-S_5-S_6-S_7\\
 &+S_8+S_9+S_{10}-2S_{11}-S_{12}-S_{13}-S_{\mathrm{sw}}.
\end{aligned}
\qquad(2\mid N,\ N\ge2).
\end{equation}


\subsection{The representation count}

Put
\begin{equation}\label{eq:source-bad-count}
\begin{aligned}
 g_b(e)&=\#\{(p_1,p_2)\in\mathbb P^2:
       N^b\le p_1\le\sqrt N,\ \varepsilon N<p_1p_2=e<N\},\\
 B_b&=\sum_e g_b(e)1_{N-e\in\mathbb P},\qquad
 P_2=\#\{n\in\mathcal A_N:\Omega(n)\le2\}.
\end{aligned}
\qquad(2\mid N,\ N\ge2).
\end{equation}
When \(e\) has distinct prime factors and both choices satisfy the
first-factor cutoff, both orientations are counted. For \(e=p^2\), there
is only the ordered pair \((p,p)\).
Hence
\begin{equation}\label{eq:source-multiplicity}
 0\le g_b(e)\le2,\qquad P_2-G_b(N)\le B_b
 \qquad(\varepsilon N>1).
\qquad(2\mid N,\ N\ge2).
\end{equation}
For the second inequality, write \(n=rs\), \(r\le s\), with both
factors prime.
If \(n\) is excluded from \(G_b\), then
\begin{equation}\label{eq:source-excluded-factor}
 N^b\le r\le\sqrt n<\sqrt N,\qquad g_b(n)\ge1.
\qquad(2\mid N,\ N\ge2,\ \varepsilon N>1).
\end{equation}
Prime squares are counted with their actual multiplicity.

Define
\begin{equation}\label{eq:source-exceptional-set}
 \mathcal E=\{n\in\mathcal A_N:P^-(n)\ge z_s,
       \ n\text{ is not squarefree or }(n,N)>1\}.
\qquad(2\mid N,\ N\ge2).
\end{equation}
For \(z_s=N^\alpha\ge2\),
\begin{equation}\label{eq:source-exceptional-size}
\begin{aligned}
 \#\{n\in\mathcal E:n\text{ is not squarefree}\}
 &\le N\sum_{k\ge\lceil z_s\rceil}k^{-2}\le\frac{2N}{z_s},\\
 \#\{n\le N:P^-(n)\ge z_s,\ (n,N)>1\}
 &\le\sum_{\substack{\ell\in\mathbb P,\ \ell\mid N\\\ell\ge z_s}}\frac N\ell
 \le\frac{11N}{z_s},\\
 |\mathcal E|&\le13N^{1-\alpha}.
\end{aligned}
\qquad(2\mid N,\ N\ge2^{1/\alpha}).
\end{equation}
The prime divisors in the middle sum are distinct.
Their number is at most
\(\lfloor1/\alpha\rfloor=11\).

Set
$N_{\mathrm{pt}}=1+\left\lceil\max\{2^{1/\alpha},
\varepsilon^{-3},\varepsilon^{-1/(1-\lambda-\beta)}\}\right\rceil$.
\begin{lemma}[Pointwise source weight bound]
\label{lem:source-pointwise}
For all $N\ge N_{\mathrm{pt}}$, the weight bounds are
\begin{equation}\label{eq:source-pointwise}
\begin{gathered}
 w(n)=0\quad(P^-(n)<z_s),\qquad
 w(n)\le4\,1_{\Omega(n)\le2}\quad(n\in\mathcal A_N\setminus\mathcal E),\\
 \max\{0,w(n)-4\,1_{\Omega(n)\le2}\}\le169
 \quad(n\in\mathcal E),
 \qquad E_{\mathrm{pt}}=0.
\end{gathered}
\qquad(2\mid N,\ N\ge N_{\mathrm{pt}}).
\end{equation}
\begin{proof}
Start from the roughness observation above.
If $\Omega(n)\ge3$, then $w(n)=\min\{W_{13}(n),0\}\le0$.
If $(n,N)>1$, the coprimality conditions force every row occurrence to vanish.
For $\Omega(n)\le2$ and $n\in\mathcal E$ with $(n,N)=1$,
only $n=p^2$ occurs; only $S_1,S_2$ contribute, giving $w(n)=4$.
For $n\notin\mathcal E$ with $\Omega(n)\le2$, a prime has weight $4$.
For $n=rs$, with primes $r<s$, the bands
$[\alpha,\beta)$, $[\beta,\mu_*)$, $[\mu_*,\lambda)$,
$[\lambda,1/3)$ and $[1/3,1/2)$ for $\log_Nr$
give weights $1,2,2,3,4$, respectively, directly from the row table.
This agrees with Paper~II, (1.11) \cite[(1.10)--(1.11)]{PaperII}.
Here $N_{\mathrm{pt}}$ is chosen so that
$\varepsilon N>N^{2/3}$ and
$\varepsilon N>N^{\lambda+\beta}$.
Then the larger prime exceeds $N^{1/3}$. The two-prime divisor
rows vanish, since their products are below $N^{\lambda+\beta}$ for
$S_8$--$S_{10}$ and below $N^{2/3}$ for $S_{11}$--$S_{13}$.
These observations prove all three bounds, including the stronger
exceptional bound with right-hand side zero.
\end{proof}
\end{lemma}
Put
\begin{equation}\label{eq:source-pointwise-error}
 E_{\mathrm{pt}}=\frac14\sum_{n\in\mathcal E}
       \max\{0,w(n)-4\,1_{\Omega(n)\le2}\}.
\qquad(2\mid N,\ N\ge2).
\end{equation}
The pointwise lemma gives $E_{\mathrm{pt}}=0$ for
$N\ge N_{\mathrm{pt}}$, within its stated scope.
We still record cruder bounds. Summation gives
\begin{equation}\label{eq:source-first-lower}
\begin{gathered}
 0\le E_{\mathrm{pt}}\le\frac{169}{4}|\mathcal E|
       \le\frac{2197}{4}N^{1-\alpha},\\
 \mathcal W\le4P_2+4E_{\mathrm{pt}},\qquad
 4G_b(N)\ge\mathcal W-4B_b-4E_{\mathrm{pt}}.
\end{gathered}
\qquad(2\mid N,\ N\ge N_{\mathrm{pt}}).
\end{equation}
The transfer to \(r\le q^{3/4}\) is \eqref{eq:intro-transfer}.

\subsection{Signed arithmetic sums}

Put
\begin{equation}\label{eq:source-fixed-sum}
\begin{aligned}
 R_{\mathrm{fix}}={}&3S_1-2S_3-S_4-S_5-S_7+S_8+S_9+S_{10}\\
 &-2S_{11}-S_{12}-S_{13}.
\end{aligned}
\qquad(2\mid N,\ N\ge2).
\end{equation}
Then \eqref{eq:source-weight-identity} gives
\begin{equation}\label{eq:source-finite-identity}
 \mathcal W-4B_b
 =R_{\mathrm{fix}}+S_2-S_6-S_{\mathrm{sw}}-4B_b.
\qquad(2\mid N,\ N\ge2).
\end{equation}
Define
\begin{equation}\label{eq:source-retained-definitions}
\begin{aligned}
 G_{\mathrm{rows}}&=2S_{11}+S_{12}+S_{13},\\
 \mathcal C_{\mathrm{ret}}&=3S_1+S_2-S_4-S_5-S_7+S_{10}-G_{\mathrm{rows}},\\
 \mathcal Q&=-2S_3+S_8+S_9.
\end{aligned}
\qquad(2\mid N,\ N\ge2).
\end{equation}
Coefficient comparison gives
\begin{equation}\label{eq:source-five-sums}
\begin{gathered}
 \mathcal C_{\mathrm{ret}}+\mathcal Q=R_{\mathrm{fix}}+S_2,\\
 \mathcal W-4B_b
 =\mathcal C_{\mathrm{ret}}+\mathcal Q-S_6-S_{\mathrm{sw}}-4B_b.
\end{gathered}
\qquad(2\mid N,\ N\ge2).
\end{equation}
Substitution proves both equalities coefficient by coefficient.
The coefficients \(3S_1\) and \(+S_{10}\) are fixed.
The integer supports may overlap.
These coefficient identities do not depend on later estimates.
Take the signed normalized chains \(Q_{j,k}\), their index set
\(\mathcal I\), and the stage index \(k\in\{0,\ldots,10\}\) from
Section~\ref{sec:retained-side}. See its ``original signed chains'' paragraph
and \eqref{eq:retB-telescope}. Each \(Q_{j,k}\) already includes its
signed normalized coefficient, so no extra multiplicity factor enters.
Finite telescoping gives
\begin{equation}\label{eq:source-signed-telescope}
 \sum_{j\in\mathcal I}Q_{j,0}
 =\sum_{j\in\mathcal I}Q_{j,10}
  +\sum_{j\in\mathcal I}\sum_{k=1}^{10}(Q_{j,k-1}-Q_{j,k}).
\qquad(2\mid N,\ N\ge2).
\end{equation}
This is the normalized-chain identity cited above. The arithmetic coefficient
identities above keep their unnormalized definitions.
Paper~II uses joint endpoints for
\(S_1,S_2\text{-}T_8/4,G_{\mathrm{rows}}\)
\cite[(6.4)--(6.4d), pp.~70--72]{PaperII}.
Here \(T_8\) is Paper~II's stage-eight endpoint in the \(S_2\) chain.
The notation \(S_2\text{-}T_8/4\) is a joint-chain label;
it defines no new arithmetic sum.
Their joint endpoint is kept before the differences are split.

\subsection{A finite upper decomposition}

Each original representation occurrence uses the shifted-prime set
of \(\eqref{eq:intro-count}\):
\begin{equation}\label{eq:source-A-N}
 \mathcal A_N=\{N-p:p\in\mathbb P,\ p<(1-\varepsilon)N\}.
\qquad(2\mid N,\ N\ge2).
\end{equation}
The two arithmetic sources stay separate. Take \(A_b,E_{\mathrm{top}}\)
from \eqref{eq:k2-source} and \(S_b\) from the Section~\ref{sec:k2} definition
\eqref{eq:k2-sifted-source}; \(A\) from \eqref{eq:l3-source};
\(\mathcal P_N\) from \eqref{eq:l3-constants}; and
\(b_A,P_N,S_v,S_A,E_{\mathrm{src}}\) from
\eqref{eq:l3-sieved-source}. The ordered prime coordinates
\(a_i\) and their prefix domains are those of \eqref{eq:l3-buchstab} and
the paragraph before it. Define
\begin{equation}\label{eq:source-local-sequences}
\begin{gathered}
 h_b(n)=A_b(N-n),\qquad
 \mathscr L_b=\{p\in\mathbb P:p\nmid2N\},\qquad E_b=E_{\mathrm{top}},\\
 h_{\mathrm{sw}}(n)=A(N-n),\qquad
 \mathscr L_{\mathrm{sw}}=\mathcal P_N,\qquad
 E_{\mathrm{sw}}^{\mathrm{src}}=E_{\mathrm{src}},
\end{gathered}
\qquad(2\mid N,\ N\ge2).
\end{equation}
Here \(A_b,E_{\mathrm{top}}\) are the Section~\ref{sec:k2} source objects in
\(\eqref{eq:k2-source}\). The objects \(A,\mathcal P_N,E_{\mathrm{src}}\) are
the Section~\ref{sec:l3} source objects in \(\eqref{eq:l3-source}\),
\(\eqref{eq:l3-constants}\), and \(\eqref{eq:l3-sieved-source}\). Both sequences are nonnegative.
No extra \(1_{(e,2N)=1}\) factor enters.
The script-font symbols \(\mathscr L_b\), \(\mathscr L_{\mathrm{sw}}\),
\(\mathscr L_X\), and \(\mathscr U_X\) denote the variable sieve families
and their sifted sums. The fixed prime support is written \(\mathcal P_N\).
For \(X\in\{b,\mathrm{sw}\}\), write \(h_X,\mathscr L_X,E_X\) for the
corresponding line of \(\eqref{eq:source-local-sequences}\), and put
\begin{equation}\label{eq:source-sifted-sequence}
\begin{aligned}
 P_X(y)&=\prod_{\substack{\ell<y\\\ell\in\mathscr L_X}}\ell,\\
 \mathscr U_X(d;y)&=\sum_{d\mid n}h_X(n)
              1_{(n/d,P_X(y))=1}.
\end{aligned}
\qquad(2\mid N,\ N\ge2).
\end{equation}
Take \(I_L,I_H,J,J_{\mathrm{tail}},J_t,z_s,\eta_{\mathrm m},
v_{\mathrm{RA}}\) from the Root--A definition \eqref{eq:ra-parameters}.
Take \(c_{R,K}\) from \eqref{eq:ra-coefficient}, and
\(R_L,R_{\mathrm{tail}},B_L,B_{\mathrm{tail}},B,c_{\mathrm{lo}},\Gamma\)
from \eqref{eq:ra-factories} and its preceding choice of fixed logarithmic
exponents. Take \(D_{\mathrm m}\) from \eqref{eq:ra-precision-order}.
Write \(Q_{\mathrm{RA}}\) for the level called \(Q\) in
\eqref{eq:ra-factories}. This alias is distinct from every signed chain
\(Q_{j,k}\). We do not redefine these objects here.
Both source comparisons require the finite threshold \(N_6\)
of Section~\ref{sec:notation}.
It enters the operative \(T_{\mathrm{RA}}\), the maximum of the
sixth-sum comparison threshold and \(N_6\). We do not redefine it here.

Use the common choice \(\sigma=\sigma_*\) from
\eqref{eq:not-sigma-choice}. It lies in the Section~\ref{sec:k2} range of
\eqref{eq:k2-parameters}, and the switched comparison uses it. Take
\(\tau,\rho\) from Section~\ref{sec:k2} \eqref{eq:k2-parameters}; Section~\ref{sec:l3} uses the same
values. We make no new choice or definition of
\(\sigma,\tau,\rho\) here.
Partition the logarithmic prime ranges into half-open intervals
\([t_i,\rho t_i)\), with \(t_i=u\rho^{k_i}\), \(k_i\ge0\).
Clip the intervals at \(\mu_*\).
Each prime has one cell.
For \(X=b\), use the retained cells of the Section~\ref{sec:k2}
\(\eqref{eq:k2-retention}\) retained-cell definition and the cutoff in
\(\eqref{eq:k2-four-rows}\). Write them as
\(\mathcal R_b,z_{b,c},\chi_{b,c}\). For \(X=\mathrm{sw}\), use the
Section~\ref{sec:l3} definitions \(\eqref{eq:l3-buchstab-peel}\) and
\(\eqref{eq:l3-row3-level}\). Write them as
\(\mathcal R_{\mathrm{sw}},z_{\mathrm{sw},c},\chi_{\mathrm{sw},c}\).
These definitions come from the two source sections.
They form no common envelope in this section.
We use no equality between the retained collections of the three sections.
The closed-cell infima of Sections~\ref{sec:k2} and~\ref{sec:l3} and their endpoint rules
apply only within their own source.
For \(X=b\), the packing function is \(H\), as specified in
\eqref{eq:k2-retention}. For \(X=\mathrm{sw}\), it is
\(\widetilde H\), as specified in \eqref{eq:l3-row3-level} and its
preceding definition of \(\Theta_c\). The two functions are defined in
Paper~II, equation (2.17d), and the paragraph after it.
We have \(\widetilde H=H\) only on the stated largest-coordinate
range, and use no equality outside it.
For either source write \(z_{X,c}=N^{t_3}\), with \(t_3\) the lower
endpoint of the third coordinate of its cell \(c\). Here
\(z_{\mathrm{sw},c}\) is the object called \(z_c\) in Section~\ref{sec:l3}; we use only the
source-indexed name. The Section~\ref{sec:k2} cutoff is defined
in \eqref{eq:k2-four-rows}, with the lower-endpoint convention stated
before \eqref{eq:k2-parent-functional}.

The source index here is always \(X\in\{b,\mathrm{sw}\}\).
For the $b$-source the underlying prime support is \(\mathcal P_N\), so
\(P_b(y)=P_N(y)\). Both sieve products use primes below $y$ not dividing
$2N$.
For the \(b\)-source, the substitution \(e=N-n\) identifies the
generic sieve with \(S_b\) of Section~\ref{sec:k2} on that source:
\begin{equation}\label{eq:source-ch07-sieve-bridge}
 \mathscr U_b(d;y)
 =\sum_{d\mid n}A_b(N-n)1_{(n/d,P_b(y))=1}
 =\sum_{d\mid N-e}A_b(e)1_{((N-e)/d,P_b(y))=1}
 =S_b(d;y).
\qquad(2\mid N,\ N\ge2).
\end{equation}
For \(X=\mathrm{sw}\), identify the generic cofactor sieve with the
Section~\ref{sec:l3} sieve on precisely the following rows:
\begin{itemize}
\item \(d=1\) with cutoff \(y=N^{\mu_*}\), the source-comparison row
\(\mathscr U_{\mathrm{sw}}(1;N^{\mu_*})=S_A\) from Section~\ref{sec:l3}
\eqref{eq:l3-sieved-source};
\item \(v=a_1\cdots a_j\), \(0\le j\le2\), with cutoff \(y=N^u\);
\item \(v=a_1a_2a_3\), with dynamic cutoff \(y=a_3\), or retained-cell
cutoff \(y=z_{\mathrm{sw},c}=N^{t_3}\);
\item \(vr=a_1a_2a_3r\) on a retained cell, with
\(z_{\mathrm{sw},c}\le r<a_3\), at cutoff \(y=r\) or
\(y=z_{\mathrm{sw},c}\).
\end{itemize}
Here the empty prefix is \(v=1\), and all primes belong to
\(\mathcal P_N\) and have the ordering in \eqref{eq:l3-buchstab}
and \eqref{eq:l3-peel}. Let \(d\) be the displayed prefix \(v\) or
product \(vr\). Every prime divisor of \(d\) is at least \(y\), so
\((d,P_{\mathrm{sw}}(y))=1\). The definitions give
\(P_{\mathrm{sw}}(y)=P_N(y)\). Write \(S_d(y)\) for the \(S_v(y)\) of Section~\ref{sec:l3}
from \eqref{eq:l3-sieved-source} with \(v=d\). For \(d\mid n\), the Section~\ref{sec:l3} definition
\eqref{eq:l3-sieved-source} gives
\begin{equation}\label{eq:source-ch08-sieve-bridge}
 \mathscr U_{\mathrm{sw}}(d;y)
 =\sum_{d\mid n}h_{\mathrm{sw}}(n)1_{(n/d,P_N(y))=1}
 =\sum_{d\mid n}b_A(n)1_{(n,P_{\mathrm{sw}}(y))=1}
 =S_d(y).
\qquad(2\mid N,\ N\ge2).
\end{equation}
This identity holds on these rows only, not for arbitrary \(d,y\).
These \(E\)-terms are unnormalized source-restoration costs. The two
source upper comparisons are stated separately:
Here \(B_b'\) is the prime-occurrence part of \(B_b\), and
\(E_{\mathrm{top}}\) is the removed-top cost.
\begin{lemma}[Upper-source comparisons]
\label{lem:source-upper-sources}
\begin{equation}\label{eq:source-upper-comparison}
\begin{gathered}
 B_b\le E_b+\mathscr U_b(1;N^{\mu_*}),\qquad
 S_{\mathrm{sw}}\le E_{\mathrm{sw}}^{\mathrm{src}}+
 \mathscr U_{\mathrm{sw}}(1;N^{\mu_*})\\
 \text{for even }N\ge N_6.
\end{gathered}
\end{equation}
\begin{proof}
The first inequality follows from Section~\ref{sec:k2} \eqref{eq:k2-source-comparison}
and \(E_b=E_{\mathrm{top}}\); the second is the Section~\ref{sec:l3} source comparison
\eqref{eq:l3-source-id}.
We do not identify their hypotheses or their restoration costs.
\end{proof}
We invoke the source comparison once, with the named hypotheses displayed above.
\end{lemma}
For real cutoffs \(1\le y\le w_{\mathrm{cut}}\), the least allowed prime gives
\begin{equation}\label{eq:source-buchstab-one-step}
 \mathscr U_X(d;w_{\mathrm{cut}})=\mathscr U_X(d;y)
 -\sum_{\substack{y\le r<w_{\mathrm{cut}}\\r\in\mathscr L_X}}\mathscr U_X(dr;r).
\qquad(2\mid N,\ N\ge2).
\end{equation}
Write \(n=dm\).
If \(m\) survives at \(y\) but fails at \(w_{\mathrm{cut}}\),
choose its least allowed prime in \([y,w_{\mathrm{cut}})\).
Exactly one summand records that prime.
Repeated copies remain in the cofactor.
Thus the identity also counts prime powers correctly.

Three applications give
\begin{equation}\label{eq:source-buchstab}
\begin{aligned}
 \mathscr U_X(1;N^{\mu_*})&=T_{X,0}-T_{X,1}+T_{X,2}-T_{X,3},\\
 T_{X,j}&=\sum_{\substack{N^u\le a_j<\cdots<a_1<N^{\mu_*}\\
                         a_i\in\mathscr L_X}}
                         \mathscr U_X(a_1\cdots a_j;N^u)\quad(0\le j\le2),\\
 T_{X,3}&=\sum_{\substack{N^u\le a_3<a_2<a_1<N^{\mu_*}\\
                         a_i\in\mathscr L_X}}
                         \mathscr U_X(a_1a_2a_3;a_3).
\end{aligned}
\qquad(2\mid N,\ N\ge2).
\end{equation}
For \(j=0\), the sum is \(\mathscr U_X(1;N^u)\).
The primes are ordered; no factorial occurs.
Every multiplicity remains in \(h_X\).

For \(d=a_1a_2a_3\) in a retained cell of
source \(X\), put
\begin{equation}\label{eq:source-dynamic-cutoff}
\begin{aligned}
 -\mathscr U_X(d;a_3)
 &= -\mathscr U_X(d;z_{X,c})+
       \sum_{\substack{z_{X,c}\le r<a_3\\r\in\mathscr L_X}}\mathscr U_X(dr;r)\\
 &\le-\mathscr U_X(d;z_{X,c})+
       \sum_{\substack{z_{X,c}\le r<a_3\\r\in\mathscr L_X}}\mathscr U_X(dr;z_{X,c}).
\end{aligned}
\qquad(2\mid N,\ N\ge2).
\end{equation}
The equality is \eqref{eq:source-buchstab-one-step}.
Nonnegativity proves the inequality.
Actual cell membership gives \(z_{X,c}\le a_3\).
The closed infimum creates no additional occurrences.

Let \(\sum_{\mathrm{ret},X}\) restrict the ordered triples to
\(\mathcal R_X\), equivalently to \(\chi_{X,c}=1\). For either source define
the peel terms
\begin{equation}\label{eq:source-base-peel}
\begin{aligned}
 R_{X,3}&=\sum_{\mathrm{ret},X}\mathscr U_X(a_1a_2a_3;z_{X,c}),\\
 P_{X,3}&=\sum_{\mathrm{ret},X}
       \sum_{\substack{z_{X,c}\le r<a_3\\r\in\mathscr L_X}}
             \mathscr U_X(a_1a_2a_3r;z_{X,c}).
\end{aligned}
\qquad(2\mid N,\ N\ge2).
\end{equation}
For each source \(X\in\{b,\mathrm{sw}\}\), the non-retained depth-three
terms are nonnegative and occur with a minus sign, so we drop them.
The retained terms use the dynamic inequality above:
\begin{equation}\label{eq:source-nonretained-drop}
 -T_{X,3}\le -R_{X,3}+P_{X,3}.
\qquad(2\mid N,\ N\ge2).
\end{equation}
For even \(N\ge N_6\), define the two upper expressions and their
slacks, one per source. The next comparison proves their nonnegativity:
\begin{equation}\label{eq:source-upper-slack}
 \begin{aligned}
 U_b&:=T_{b,0}-T_{b,1}+T_{b,2}-R_{b,3}+P_{b,3},&
 Z_b&:=E_b+U_b-B_b\ge0,\\
 U_{\mathrm{sw}}&:=T_{\mathrm{sw},0}-T_{\mathrm{sw},1}
       +T_{\mathrm{sw},2}-R_{\mathrm{sw},3}+P_{\mathrm{sw},3},&
 Z_{\mathrm{sw}}&:=E_{\mathrm{sw}}^{\mathrm{src}}
       +U_{\mathrm{sw}}-S_{\mathrm{sw}}\ge0.
 \end{aligned}
 \qquad(2\mid N,\ N\ge N_6).
\end{equation}
Hence
\begin{equation}\label{eq:source-sifted-upper}
 \mathscr U_X(1;N^{\mu_*})\le
 \begin{cases}
 U_b,&X=b,\\
 U_{\mathrm{sw}},&X=\mathrm{sw}.
 \end{cases}
\qquad(2\mid N,\ N\ge N_6).
\end{equation}
These inequalities assume the two source comparisons in
\eqref{eq:source-upper-comparison}. They give no estimate for either
restoration cost.

\subsection{The conditional lower expression}

Take \(A_{\mathrm{mix}}\) from the Root--A definition
\eqref{eq:ra-mix-definition}, and \(Z_6\) with its comparison from
Section~\ref{sec:root-a} \eqref{eq:ra-positive-comparison}. We redefine neither object.
This sum uses the ordinary upper Rosser weights
\(\mu_d^+(D)\) with the odd-index cubic gates of
\(\eqref{eq:ra-rosser-weight}\). Their upper-sieve property is the
Iwaniec upper-bound lemma
\cite[Lemma~1, (3.2), (3.6)]{Iwaniec1980Rosser}
\cite[Lemmas~2--3, p.~313]{Iwaniec1980Remainder}. This section
introduces no new weight.
The Section~\ref{sec:root-a} comparison gives
\begin{equation}\label{eq:source-mixed-comparison}
 S_6\le A_{\mathrm{mix}},\qquad Z_6\ge0
 \qquad(2\mid N,\ N\ge N_6).
\end{equation}
For a mark \(p\), put \(m=(N-e)/p\) and
\(P_{\mathrm{RA}}=\prod_{\ell<z_s,\,\ell\nmid2N}\ell\).
For \(D>1\), the cited Iwaniec upper-bound lemmas give, for the upper
Rosser divisor sum,
\cite[Lemma~1, (3.2), (3.6)]{Iwaniec1980Rosser}
\cite[Lemmas~2--3, p.~313]{Iwaniec1980Remainder},
\[
 \sum_{d\mid(m,P_{\mathrm{RA}})}\mu_d^+(D)
 \ge1_{(m,P_{\mathrm{RA}})=1}\ge1_{P^-(m)\ge z_s}
 \qquad(2\mid N,\ N\ge2,\ D>1).
\]
Dropping \((m,Np)=1\) enlarges the original nonnegative occurrence count.
For even \(N\ge N_6\), each of \(R_L,R_{\mathrm{tail}},Q_{\mathrm{RA}}\) exceeds
\(N^{1/3}\). Hence take \(D=R_L/p,R_{\mathrm{tail}}/p,Q_{\mathrm{RA}}/p\) on,
respectively, \(I_L\times J\), \(I_H\times J_{\mathrm{tail}}\),
\(I_H\times J_t\). These pieces partition the original source: \(e=N/4\)
belongs to \(I_L\), and the shared mark endpoint belongs to
\(J_{\mathrm{tail}}\).
Since \(p\ge N^{t_0}>z_s\), the mark is unique in \(q=pd\).
The condition \(d\mid m\) is exactly \(e\equiv N\pmod{pd}\). The original gates give
\(c_{R,K}(pd)=\mu_d^+(R/p)\). Summation gives precisely
\(S_6\le A_{\mathrm{mix}}\) and \(Z_6\ge0\), with all mark
multiplicities. A product-cell partition only splits these exact
sums. It must keep the original \(pd<R\), the cubic gates and every clipped
endpoint. In particular \(c_{\mathrm{lo}}\) and \(\Gamma\) keep the
gates at \(Q_{\mathrm{RA}}\).
Subdivide each of the three marked pieces in
\eqref{eq:ra-three-pieces} into cells in its two displayed coordinates.
Let \(\delta=L^{-D_{\mathrm m}}\), and use half-open multiplicative
intervals of relative width at most \(1+\delta\) in each coordinate.
Insert all row cuts, the points
\(N/4\), \(1/3-\eta_{\mathrm m}\), \(1/3\), and the two endpoints of
\(I_H\) as cell boundaries. Denote the resulting finite, disjoint cell
 families by \(\Theta^{\sharp}_L,\Theta^{\sharp}_{\mathrm{tail}},\Theta^{\sharp}_t\).
Use the separate piece index \(Y\in\{L,\mathrm{tail},t\}\) and the
mapping
\begin{equation}\label{eq:source-root-a-piece-map}
 (I_Y,K_Y,R_Y)=
 \begin{cases}
 (I_L,J,R_L),&Y=L,\\
 (I_H,J_{\mathrm{tail}},R_{\mathrm{tail}}),&Y=\mathrm{tail},\\
 (I_H,J_t,Q_{\mathrm{RA}}),&Y=t.
 \end{cases}
\qquad(2\mid N,\ N\ge N_6).
\end{equation}
Thus each \(C\in\Theta^{\sharp}_Y\) satisfies
 \(C\subseteq(\mathbb P\cap I_Y)\times(\mathbb P\cap N^{K_Y})\) and
\((p,N)=1\), with \(0<(N-e)/p\) on every included occurrence. Their unions
are exactly the corresponding coprime marked source pieces
\(I_L\times J\), \(I_H\times J_{\mathrm{tail}}\), and \(I_H\times J_t\).
A cell touching a boundary is kept as a boundary cell.
For \(Y\in\{L,\mathrm{tail},t\}\) and \(C\in\Theta^{\sharp}_Y\), define
\begin{equation}\label{eq:source-cell-sums}
\begin{aligned}
 S_{6,Y,C}&:=\sum_{(e,p)\in C}\mathbf1_{p\in\mathbb P}\mathbf1_{(p,N)=1}\mathbf1_{p\mid N-e}\mathbf1_{0<(N-e)/p}\mathbf1_{e\in\mathbb P}
 \mathbf1_{P^-((N-e)/p)\ge z_s}\mathbf1_{((N-e)/p,Np)=1},\\
 A_{Y,C}&:=\sum_{(e,p)\in C}\mathbf1_{p\in\mathbb P}\mathbf1_{(p,N)=1}\mathbf1_{p\mid N-e}\mathbf1_{0<(N-e)/p}\mathbf1_{e\in\mathbb P}
 \sum_{\substack{d\mid (N-e)/p\\d\mid P_{\mathrm{RA}}}}
 c_{R_Y,K_Y}(pd)\mathbf1_{e\equiv N\ (pd)},
\end{aligned}
\qquad(2\mid N,\ N\ge2).
\end{equation}
Here \(R_Y\) is \(R_L\), \(R_{\mathrm{tail}}\), or \(Q_{\mathrm{RA}}\). The
corresponding coefficient gates stay in the inner sum. The
half-open families make every marked occurrence belong to exactly one
\(C\), including \(e=N/4\) and the shared mark endpoint. The one-term
Rosser inequality at each occurrence gives \(S_{6,Y,C}\le A_{Y,C}\), cell
by cell. Reindexing the finite sums gives
\begin{equation}\label{eq:source-cell-sum}
 S_6=\sum_Y\sum_{C\in\Theta^{\sharp}_Y}S_{6,Y,C},\qquad
 \sum_Y\sum_{C\in\Theta^{\sharp}_Y}A_{Y,C}=A_{\mathrm{mix}}.
\qquad(2\mid N,\ N\ge N_6).
\end{equation}
The finite comparison uses these cellwise Rosser inequalities.
Thus the displayed comparison is the sum of finitely many explicit cell
inequalities. No product cell, endpoint, or boundary strip is
omitted. The later complete-cell estimates and boundary-strip bound are
quantitative inputs only and do not alter this finite identity.
The three slacks below are cited or defined by the source comparisons
above, before this substitution.
Substitution in \eqref{eq:source-finite-identity} gives
\begin{equation}\label{eq:source-exact-slack}
\begin{aligned}
 \mathcal W-4B_b-4E_{\mathrm{pt}}
 ={}&R_{\mathrm{fix}}+S_2-A_{\mathrm{mix}}
       -U_{\mathrm{sw}}-4U_b\\
    &-E_{\mathrm{sw}}^{\mathrm{src}}-4E_b-4E_{\mathrm{pt}}
       +Z_6+Z_{\mathrm{sw}}+4Z_b.
\end{aligned}
\qquad(2\mid N,\ N\ge N_6).
\end{equation}
This equality follows from the source identity and the cited and
local definitions. The Section~\ref{sec:root-a} comparison
\eqref{eq:ra-positive-comparison} gives \(Z_6\ge0\). The comparisons in
\eqref{eq:source-upper-slack} give \(Z_{\mathrm{sw}},Z_b\ge0\).
For even \(N\ge\max\{N_6,N_{\mathrm{pt}}\}\),
\eqref{eq:source-first-lower} and nonnegativity of these slacks yield
\begin{equation}\label{eq:source-final-lower}
\begin{aligned}
 4G_b(N)\ge{}&R_{\mathrm{fix}}+S_2-A_{\mathrm{mix}}\\
 &-(T_{\mathrm{sw},0}-T_{\mathrm{sw},1}+T_{\mathrm{sw},2}
                         -R_{\mathrm{sw},3}+P_{\mathrm{sw},3})\\
 &-4(T_{b,0}-T_{b,1}+T_{b,2}-R_{b,3}+P_{b,3})\\
 &-E_{\mathrm{sw}}^{\mathrm{src}}-4E_b-4E_{\mathrm{pt}}.
\end{aligned}
\qquad(2\mid N,\ N\ge\max\{N_6,N_{\mathrm{pt}}\}).
\end{equation}
The two expanded signs are
\begin{equation}\label{eq:source-expanded-signs}
\begin{aligned}
 -U_{\mathrm{sw}}&=-T_{\mathrm{sw},0}+T_{\mathrm{sw},1}
                  -T_{\mathrm{sw},2}+R_{\mathrm{sw},3}-P_{\mathrm{sw},3},\\
 -4U_b&=-4T_{b,0}+4T_{b,1}-4T_{b,2}+4R_{b,3}-4P_{b,3}.
\end{aligned}
\qquad(2\mid N,\ N\ge N_6).
\end{equation}
Use one bound for the complete \(S_{\mathrm{sw}}\).
Paper~II gives its coefficient-eight estimate
\cite[(1.9), (5.29)--(5.43)]{PaperII}; the Buchstab upper bound
concerns the same occurrence.
Sections~\ref{sec:k2} and \ref{sec:l3} estimate these two five-term
upper expressions (four principal rows and their peel term) separately.
The displayed expansion fixes each peel term's sign.
The finite switched-source domination is already in the
local \(Z_{\mathrm{sw}}\) identity above. The remaining distribution
and quantitative estimates for this occurrence are separate hypotheses
of the later bill. They lie outside the finite decomposition of this section.


\subsection{\texorpdfstring{The lower expression with the original $S_2$}{The lower expression with the original S2}}

Keep the original \(S_2\) throughout.
Use $\overline H$, $M_{5,\mathrm{internal}}$ and $\Delta D_{L3}$
from \eqref{eq:intro-components}.
Here $\overline H$ is the switched integral bound.
The distinct constant $153/256$ bounds the packing envelope.
The complete normalized switched error is $E_{\mathrm{sw}}(N;\sigma)$.
The source-restoration error $E_{\mathrm{sw}}^{\mathrm{src}}$ is unnormalized.
The main route takes $g_6:=0$ and excludes the optional six-endpoint
improvement.
Fix this common parameter and an even \(N\ge\max\{N_6,N_{\mathrm{pt}}\}\)
throughout the following conditional lemma. Take
\(W,E_{\mathrm{van}},\Lambda^0_{j,N},\Lambda_N,\Delta_j,\Delta_{G1}\)
from \eqref{eq:not-normalization}, \eqref{eq:intro-retained-error}, and
Section~\ref{sec:retained-side}, and \(g_{\mathrm{RA,lo}},E_{\mathrm{RA}}\)
from \eqref{eq:ra-b11-floor}. Use the complete switched and exterior
errors from \eqref{eq:l3-final} and \eqref{eq:k2-final-fixed-sigma}.
The fixed constants are defined in \eqref{eq:intro-components}
and \eqref{eq:intro-fallback-constant}.

\begin{lemma}[Addition of the four comparisons]
\label{lem:source-trackb-combination}
Assume \eqref{eq:source-pointwise} and the following four inequalities.
Write
\begin{equation}\label{eq:source-trackb-retained}
\begin{aligned}
 \frac{R_{\mathrm{fix}}+S_2}{W(N)}
 &\ge D_{\mathrm{II}}+\overline J_6
       +2\overline H+\frac{5719}{10^6}
       +M_{5,\mathrm{internal}}-E'(N;\sigma),\\
 E'(N;\sigma)&=E_{\mathrm{van}}(N;\sigma)
    +\sum_{j\in\mathcal H}
          \bigl(|\Lambda^0_{j,N}|+|\Lambda_N(\Delta_j)|\bigr)
    +|\Lambda_N(\Delta_{G1})|.
\end{aligned}
\qquad(2\mid N,\ N\ge\max\{N_6,N_{\mathrm{pt}}\}\ \text{(assumed)}).
\end{equation}
Here \(D_{\mathrm{II}}\) denotes the original Paper~II floor.
The main route does not use the six endpoint replacements.
Here $\mathcal H=\{4,5,7,10\}$, with $5$ denoting the lower
part of $S_5$ and $10$ its retained $S_{10}$ part.
Section~\ref{sec:retained-side} gives the hypotheses for this estimate.

The other three comparisons have the forms
\begin{equation}\label{eq:source-trackb-three-estimates}
\begin{aligned}
 \frac{S_6}{W(N)}
   &\le\overline J_6-g_{\mathrm{RA,lo}}+E_{\mathrm{RA}}(N;\sigma),\\
 \frac{S_{\mathrm{sw}}}{W(N)}
   &\le2\overline H-\Delta D_{L3}+E_{\mathrm{sw}}(N;\sigma),\\
 \frac{4B_b}{W(N)}&\le c_{K2,G0}I(b)+E_K(N;b,\sigma).
\end{aligned}
\qquad(2\mid N,\ N\ge\max\{N_6,N_{\mathrm{pt}}\}\ \text{(assumed)}).
\end{equation}
All error terms are nonnegative.

Then the following normalized lower bound holds:
\begin{equation}\label{eq:source-trackb-combination}
\begin{aligned}
 \frac{4G_b(N)}{W(N)}
 &\ge D_{\mathrm{II}}+\frac{5719}{10^6}
      +M_{5,\mathrm{internal}}+g_{\mathrm{RA,lo}}+\Delta D_{L3}
      -c_{K2,G0}I(b)-\mathcal E_B(N;b,\sigma),\\
 \mathcal E_B(N;b,\sigma)
 &=E'(N;\sigma)+E_{\mathrm{RA}}(N;\sigma)+E_{\mathrm{sw}}(N;\sigma)+E_K(N;b,\sigma)
       +\frac{4E_{\mathrm{pt}}}{W(N)}.
\end{aligned}
\qquad(2\mid N,\ N\ge\max\{N_6,N_{\mathrm{pt}}\}).
\end{equation}
\end{lemma}

\begin{proof}
Substitute the four assumed estimates into
\eqref{eq:source-finite-identity} and use \eqref{eq:source-first-lower}.
The terms \(\overline J_6\) and \(2\overline H\) cancel exactly.
Each complete upper source is subtracted once.
In particular, no additional switched majorant is added.

Using the same retained constant as in Section~\ref{sec:final-bill}, put
\begin{equation}\label{eq:source-trackb-floor}
\begin{aligned}
 D_B&=D_{\mathrm{II}}+\frac{5719}{10^6}
      +M_{5,\mathrm{internal}}+g_{\mathrm{RA,lo}}
      -\frac1{10000}+\Delta D_{L3},\\
 \mathcal E_B(N;b,\sigma)\le\frac1{10000}
 &\quad\Longrightarrow\quad
 \frac{4G_b(N)}{W(N)}\ge D_B-c_{K2,G0}I(b).
\end{aligned}
\qquad(2\mid N,\ N\ge\max\{N_6,N_{\mathrm{pt}}\}).
\end{equation}
\end{proof}
